\documentclass[10pt,a4paper]{article}
\input{Formato_2021}
%--------------------------------------------------------------------------------------------------
\begin{document}

% PORTADA %%%%%%%%%%%%%%%%%%%%%%%%%%%%%%%%%%%%%%%%%%%%%%%%%%%%%%%%%%%%%%%%%%%%%%%%%%%%%%%%%%%%%%%%%%
\pagestyle{empty}
\vspace{2cm}
\begin{center}
\includegraphics[width=0.7\linewidth]{Imagenes/logo_EPEF.jpg}
\end{center}
\vspace{0.8cm}
\begin{center}
\includegraphics[width=0.7\linewidth]{Imagenes/logo_UDCb.png}
% COMENTAR LINEA ANTERIOR Y DESCOMENTAR SIGUIENTE SI SE QUIERE LOGO DE LA UDC A COLOR
%\includegraphics[width=0.7\linewidth]{Imagenes/logo_UDC.png}
\end{center}

\large

\vspace{3cm}

\begin{center}
{\setlength\arrayrulewidth{2pt}
\begin{tabular}{r|p{9.8cm}}
\arrayrulecolor{udc}
\colorbox{udc}{\textcolor{white}{\bf TÍTULO}}      
&	TÍTULO DEL TRABAJO (Si es largo, se distribuirá automáticamente en varias líneas, cuya anchura está establecida en el formato)         \\[2cm]
%\colorbox{udc}{\textcolor{white}{\bf GRADO}}       & Dual en Ingeniería Eléctrica    \\[1cm]
\colorbox{udc}{\textcolor{white}{\bf GRADO}}       & Ingeniería Electrónica Industrial y Automática    \\[1cm]
\colorbox{udc}{\textcolor{white}{\bf ASIGNATURA}}  & Nombre de la asignatura     \\[2cm]
\colorbox{udc}{\textcolor{white}{\bf ESTUDIANTE}}  &	Apellido1 Apellido 2, Nombre  \\[2cm]
%                                                  &	Apellido1 Apellido 2, Nombre  \\[2cm]
\colorbox{udc}{\textcolor{white}{FECHA}}       &	mes de 20XX
\end{tabular}}
\end{center}
\normalsize
% FIN DE LA PORTADA %%%%%%%%%%%%%%%%%%%%%%%%%%%%%%%%%%%%%%%%%%%%%%%%%%%%%%%%%%%%%%%%%%%%%%%%%%%%%%%%%%%%%%%%%%%%%%%%%%%%%%%%%%%%% -----------------------------------------------------------------------------------------------------------------------------------
\cleardoublepage

% Índice de contenidos
\pagestyle{plain}
\renewcommand{\contentsname}{Índice}
{\hypersetup{hidelinks}\tableofcontents}

\cleardoublepage

% Página que contiene el índice de listas de figuras
\phantomsection
\renewcommand*\listfigurename{Listado de figuras}
\addcontentsline{toc}{section}{\listfigurename}
{\hypersetup{hidelinks}\listoffigures}
\addtocontents{lof}{\protect\thispagestyle{plain}}

% Página que contiene el índice de listas de tablas
%\cleardoublepage
\phantomsection
\renewcommand*\listtablename{Listado de tablas}
\addcontentsline{toc}{section}{\listtablename}
{\hypersetup{hidelinks}\listoftables}
\addtocontents{lot}{\protect\thispagestyle{plain}}


% Página que contiene el índice de listados de programación
\renewcommand\lstlistlistingname{Listado de códigos de programación}
\renewcommand\lstlistingname{Código}
%\cleardoublepage
\phantomsection
\addcontentsline{toc}{section}{\lstlistlistingname}
{\hypersetup{hidelinks}\lstlistoflistings}
{\protect\thispagestyle{plain}}

% Final del índice -----------------------------------------------------------------------------------------------------------------------------------

\cleardoublepage %-----------------------------------------------------------------------------------------------------

\pagestyle{fancy}

\addcontentsline{toc}{section}{Introducción}%--------------------------------------------------------------------------------
\section*{Introducción}


Esta sección debe presentar un breve resumen de los contenidos del trabajo. Se comenta el desarrollo de las restantes secciones del trabajo, destacando las más importantes. Su extensión debe ser acorde a la del trabajo completo. Como norma general podemos indicar que su extensión no debe superar la décima parte del total. 
\href{https://www.youtube.com/watch?v=Mpf5CIkv7TY&pp=ygUUZWwgbWFtYm8ga2lrbyByaXZlcmE%3D}{UDC}

\section{Primera Sección}
\subsection{Subsección uno}%------------------------------------------------------------------------------------------------------


Contenido de la primera subsección. Se incluye como ejemplo la Figura (\ref{fig:G}).

\begin{figure}[H]
\centering
\includegraphics[scale=0.15]{Imagenes/graduacion.jpg}
\caption[Graduación en la EUP. Fuente: https://www.udc.es]{Graduación en la EUP}
\label{fig:G}
\end{figure}


A continuación, la Figura (\ref{fig:EP}).

\begin{figure}[H]
\centering
\includegraphics[scale=0.5]{Imagenes/EP.PNG}
\caption[Conjunto. Fuente: elaboración propia]{Conjunto}
\label{fig:EP}
\end{figure}


\begin{figure}[H]
\centering
\subfigure[Tipos de circuito]
{\centering
\includegraphics[width=4cm]{Imagenes/Tipos_circuito.png}
\label{TC}}
\hspace{1cm}
\subfigure[Elementos]
{\centering
\includegraphics[width=5cm]{Imagenes/elementos_circuito.jpg}
\label{EC}}
\caption[Circuitos eléctricos. Fuente: https://electromundo.pro]{Circuitos eléctricos.}\label{CC}
\end{figure}

La Figura (\ref{CC}) tiene dos elementos. La Subfigura (\ref{TC}) contiene tipos de circuitos, y la Subfigura (\ref{EC}) contiene elementos.



\subsection{Subsección dos}%------------------------------------------------------------------------------------------------------

Contenido de la segunda subsección. Se incluyen dos ejemplos de tablas.

\begin{table}[H]
\centering
\begin{tabular}{cc}
País & Capital \\ \hline
España   & Madrid \\  
Francia  & París \\  
Alemania & Berlín \\ \hline
\end{tabular}
\caption[Países y capitales. Fuente: elaboración propia]{Países y capitales}
\label{tab:PyC}
\end{table}

Esta es la Tabla \ref{tab:PyC}.


\section{Códigos de programación}%-------------------------------------------------------------------------------------------------

En esta sección se incluyen los códigos de programación. Se pueden alternar códigos escritos en distintos lenguajes de programación. Para ello se utilizará el comando \verb|\lstinputlisting[]{nombre_del_fichero}|.
Entre los corchetes se podrán introducir parámetros, separados por comas, como por ejemplo:
\begin{itemize}
    \item \textit{language}: Para indicar el lenguaje de programación del código.
    \item \textit{firstline}: Para indicar la primera línea de código a mostrar.
    \item \textit{lastline}: Para indicar la última línea de código a mostrar.
    \item \textit{firstnumber}: Para indicar desde que número se debe empezar el conteo de lineas que se verá a la izquierda del código. Normalmente, este debe ser igual a \textit{firstline}.
    \item \textit{caption}: Para añadir un titulo al código.
    \item \textit{item}: Identificador que permite realizar referencias a este código.
\end{itemize}


\bigskip


\section{Fórmulas químicas}%------------------------------------------------------------------------------------------------------

Las reacciones químicas se deben escribir como una ecuación matemática. El código \verb/\longrightarrow/ genera una fecha que separa las partes de la reacción, mientras que los símbolos \verb/~/ generan espacios en blanco. Un ejemplo de ecuación es:

\bigskip

\begin{minipage}{0.45\textwidth}
\begin{verbatim}
$${\text{Zn}}{\scriptstyle (s)}~
+~2~\text{HCl}{\scriptstyle(ac)}~
\longrightarrow 
~\text{ZnCl}_{2}{\scriptstyle (ac)}~ 
+~\text{H}_{2}{\scriptstyle (g)}$$
\end{verbatim}
\end{minipage}
\noindent\begin{minipage}{0.55\textwidth}
$${\text{Zn}}{\scriptstyle (s)}~
+~2~\text{HCl}{\scriptstyle(ac)}~
\longrightarrow 
~\text{ZnCl}_{2}{\scriptstyle (ac)}~ 
+~\text{H}_{2}{\scriptstyle (g)}$$
\end{minipage}

\bigskip

El código situado en la parte izquierda se puede copiar de este documento .pdf y pegarlo en un fichero .tex. A continuación se puede modificar para escribir otras reacciones.



\bigskip

Para escribir la identificación de un elemento químico con su masa se debe escribir el símbolo del elemento acompañado de su número atómico (identificación) como subíndice y su masa atómica como superíndice, ambos a la izquierda del símbolo del elemento. Para facilitar la escritura se ha generado un código que da formato a los datos del elemento. Este código es: 

\bigskip

\verb/                         \IE{símbolo}{número atómico}{masa atómica}/. 

\newpage

Algunos ejemplos son:

\begin{center}
\begin{tabular}{ll}
\verb/\IE{C}{12}{6}/     &  \IE{C}{12}{6}  \\ 
\verb/\IE{Ne}{22}{10}/   &  \IE{Ne}{22}{10}  \\ 
\verb/\IE{Na}{23}{11}/   &  \IE{Na}{23}{11}  \\ 
\end{tabular}
\end{center}

\section{Sección Planos}%------------------------------------------------------------------------------------------------------

A continuación se introducen los planos. Si se cambia el formato a doble cara, cada plano figurará en una página separada con numeración impar. La parte posterior de esas hojas estará en blanco.

%\pagestyle{fancy}


% FIN DE LA PORTADA
%%%%%%%%%%%%%%%%%%%%%%%%%%%%%%%%%%%%%%%%%%%%%%%%%%%%%%%%%%%%%%%%%%%%%%%%%%%%%%%%%%%%%%%%%%%%%%%%%%%%%%%%%%%%%%%%%%%%%%%%%%%%%%%%%%%%%%

% Principio Plano A4 ---------------------------------------------------------------------------------------------------
\cleardoublepage

\begin{Plano}[!htb]
\addcontentsline{toc}{subsection}{Nombre del plano A4} % Nombre que aparecerá en la lista de planos
\includepdf[pages=-]{Planos/planoA4.pdf} % Ruta y nombre del plano A4
\caption[Nombre del plano A4]{} % Página \thepage}
\label{etiquetaplanoA4} % Etiqueta para referenciar el plano
\end{Plano}
% Fin Plano A4 --------------------------------------------------------------------------------------------------------


% Principio Plano A3 ---------------------------------------------------------------------------------------------------
\cleardoublepage

\begin{landscape} % se indica que empieza una página apaisada
\afterpage

\KOMAoptions{paper=A3,paper=landscape,pagesize}
\recalctypearea
\begin{Plano}[!htb]
\addcontentsline{toc}{subsection}{Nombre del plano A3} % Nombre que aparecerá en la lista de planos
\includepdf[pages=-]{Planos/planoA3.pdf} % Ruta y nombre del plano A3
\caption[Nombre del plano A3]{}
\label{etiquetaplanoA3} % Etiqueta para referenciar el plano
\end{Plano}

\cleardoublepage

\KOMAoptions{paper=A4,pagesize}
\recalctypearea
\end{landscape} % se indica que acaba una página apaisada
% Fin Plano A3 --------------------------------------------------------------------------------------------------------


% Principio Plano A2 ---------------------------------------------------------------------------------------------------
\cleardoublepage

\begin{landscape} % se indica que empieza una página apaisada
\afterpage

\KOMAoptions{paper=A2,paper=landscape,pagesize}
\recalctypearea
\begin{Plano}[!htb]
\addcontentsline{toc}{subsection}{Nombre del plano A2} % Nombre que aparecerá en la lista de planos
\includepdf[pages=-]{Planos/planoA2.pdf} % Ruta y nombre del plano A2
\caption[Nombre del plano A2]{} % Nombre que aparecerá en la lista de planos
\label{etiquetaplanoA2} % Etiqueta para referenciar el plano
\end{Plano}

\cleardoublepage

\KOMAoptions{paper=A4,pagesize}
\recalctypearea
\end{landscape} % se indica que acaba una página apaisada
% Fin Plano A2 ---------------------------------------------------------------------------------------------------



\cleardoublepage
%%%%%%%%%%%%%%%%%%%%%%%%%%%%%%%%%%%%%%%%%%%%%%%%%%%%%%%%%%%%%%%%%%%%%%%%%%%%%%%%%%%%%%%%%%%%%%%%%%%%%%%%%%%%%%%%%%%%%%%%%%%%%%%%%%%%%%


\vspace{1cm}

% Referencias ---------------------------------------------------------------------------------------------------------

\section{Normas y referencias}



Todas las citas bibliográficas se deben indicar con una línea que comienza por \verb/\bibitem{Etiqueta}/, siendo \verb/Etiqueta/ la referencia que usaremos posteriormente para citarlas. Por ejemplo, para citar la primera norma legal, que en el fichero \LaTeX \ tiene como etiqueta {\tt UFD} se debe escribir \verb/\cite{UFD}/, apareciendo en el {\tt pdf} el valor \cite{UFD}.


\renewcommand\bibname{Disposiciones legales y normas aplicadas}


\begin{thebibliography}{X}% Si hay más de 10 referencias se debe escribir \begin{thebibliography}{XX}
% Ordenar alfabéticamante atendiendo a la primera letra
% Ejemplo de Norma UNE
\bibitem{AEN} \textsc{Asociación Española De Normalización Y Certificación}. UNE 157001: 2014. 
Criterios generales para la elaboración formal de los documentos que constituyen un proyecto técnico. Madrid: AENOR, 2014.

%Ejemplo de Legislación
\bibitem{E} \textsc{España}. Real Decreto 2267/2004, de 3 de diciembre, por el que se aprueba el Reglamento de seguridad contra incendios en los establecimientos industriales. Boletín Oficial del Estado, 17 diciembre 2004, nº. 303, págs. 41194-41255.

% Ejemplo de Norma en Internet 
\bibitem{UFD} \textsc{Unión Fenosa Distribución}. Especificaciones Particulares para Instalaciones de Medida de Energía en redes de Un $<=$ 20kV[en línea]. Madrid: Unión Fenosa Distribución, 2010 [Consulta: 1 agosto 2014]. Disponible en: 
\url{http://www.unionfenosadistribucion.com/servlet/ficheros/1297097599924/265\%5C28\%5CES_0130_ES_RE_EMA_MMT_accesible,1.pdf}
\setcounter{contador}{\value{enumiv}}
\end{thebibliography}

\renewcommand\bibname{Bibliografía}

\begin{thebibliography}{X}% Si hay más de 10 referencias se debe escribir \begin{thebibliography}{XX}
% Ordenar alfabéticamante atendiendo a la primera letra
\setcounter{enumiv}{\value{contador}}
% Ejemplo de revista en papel
\bibitem{REVISTA} \textsc{Lamas Galdo, M.I., et al.} \emph{Modelo de mecánica de fluidos computacional para el proceso de barrido en un motor Otto de dos tiempos}. (2011), 86(2), págs. 165-172.
% Ejemplo de libro en papel
\bibitem{LIBRO} \textsc{Montero, J.}, \emph{Métodos matemáticos aplicados a la ganadería.}, 3ª ed. Sevilla, Ediciones de la pradera, 2007.


\setcounter{contador}{\value{enumiv}}
\end{thebibliography}

\renewcommand\bibname{Programas de cálculo}
% Ordenar alfabéticamante atendiendo a la primera letra

\begin{thebibliography}{X}% Si hay más de 10 referencias se debe escribir \begin{thebibliography}{XX}
\setcounter{enumiv}{\value{contador}}
% Ejemplo de software en línea ARDUINO. 
\bibitem{AS} \textsc{Arduino} [software]. Versión 1.0.5. 15 Mayo 2013 [consulta: 4 agosto 2014]. 
Disponible en: \url{http://arduino.cc/en/Main/Software}.
% Ejemplo de software en CD-ROM
\bibitem{AL} \textsc{Autodesk}. AutoCAD 2012 [software], San Rafael [California]: Autodesk, 2012 [consulta: 15 agosto 2013].
\setcounter{contador}{\value{enumiv}}
\end{thebibliography}


\renewcommand\bibname{Otras referencias}

\begin{thebibliography}{X}% Si hay más de 10 referencias se debe escribir \begin{thebibliography}{XX}
% Ordenar alfabéticamante atendiendo a la primera letra
\setcounter{enumiv}{\value{contador}}
% Ejemplo de página web
\bibitem{REE} \textsc{Red Eléctrica de España} [en línea]. Madrid: REE, 2012 [consulta: 12 marzo 2012]. Disponible en: \url{http://www.ree.es/}.
%Ejemplo de catálogo de tarifas:
\bibitem{SC}\textsc{Schneider Electric}. Infraestructura de carga rápida de vehículos eléctricos [en línea]. Barcelona: Schneider Electric España, 2020 [consulta: 11 marzo 2020]. Disponible en: 
\url{https://www.se.com/es/es/product-range/60852-carga-r\%C3\%A1pida/?filter=business-4-distribuci\%C3\%B3n-el\%C3\%A9ctrica-en-baja-tensi\%C3\%B3n}.
% Ejemplo de manual de equipamiento
\bibitem{RA}\textsc{Raspberry Pi Foundation}. Raspberry Pi Documentation [en línea].  United Kingdom: Raspberry Pi, 2020 [consulta: 11 marzo 2020]. Disponible en: \url{https://www.raspberrypi.org/documentation/}.

\end{thebibliography}




\end{document}
